
\begingroup
\centering
\scriptsize
\begin{tabular}{lcc}
   \tabularnewline \midrule \midrule
   Dependent Variable: & \multicolumn{2}{c}{welfare}\\
                    & $\alpha=0$     & $\alpha=0.5$ \\    
   Model:           & (1)            & (2)\\  
   \midrule
   \emph{Variables}\\
   $\pi$ deviation  & 0.1917$^{***}$ & 0.1914$^{***}$\\   
                    & (0.0643)       & (0.0644)\\   
   treatment\_2     & -0.0048        & -0.0049\\   
                    & (0.0250)       & (0.0250)\\   
   \midrule
   \emph{Fixed-effects}\\
   segment          & Yes            & Yes\\  
   n                & Yes            & Yes\\  
   \midrule
   \emph{Fit statistics}\\
   Observations     & 255            & 255\\  
   R$^2$            & 0.27126        & 0.27112\\  
   Adjusted R$^2$   & 0.25061        & 0.25046\\  
   \midrule \midrule
   \multicolumn{3}{l}{\emph{Clustered (global\_group\_id) standard-errors in parentheses}}\\
   \multicolumn{3}{l}{\emph{Signif. Codes: ***: 0.01, **: 0.05, *: 0.1}}\\
\end{tabular}
 
\par \raggedright 
\textit{Notes:} $\pi$ deviation $= \pi_\text{sale} - \pi^*$, where $\pi_\text{sale}$ is the seller's posterior at the moment of sale and $\pi^*$ is the M\&M (2020) Appendix D equilibrium threshold for the seller's sale position $n$, risk aversion $\alpha$, and treatment-specific liquidation rule (T1 $\to$ random, T2 $\to$ average). Fixed effects: segment, sale position $n$. Standard errors clustered at the group level. We omit asymmetric (linear-spline-at-zero) and indicator specifications for the below-threshold region because only 6 sales at $\alpha=0$ and 8 at $\alpha=0.5$ fall below the threshold; at $\alpha=0$ all 6 share $\pi_\text{deviation} \approx -0.004$, so the below-threshold slope and indicator are not identified from within-sample variation. Significance codes: *** $p<0.01$, ** $p<0.05$, * $p<0.1$.
\par\endgroup


